\section{Minimal Planning Workflow}

\begin{templatebox}
\textbf{Minimal transition plan}
\begin{enumerate}[nosep]
  \item Name the target state.
  \item Name the likely decision window.
  \item Write the first visible action.
  \item Place the required objects in the future environment.
  \item Define the stop rule for the session.
\end{enumerate}
\end{templatebox}

\begin{remark}
The workflow is pedagogical rather than theorem-like. Its justification is that
each line targets one of the variables formalized above: target state fixes the
objective, decision window fixes timing, first action lowers ambiguity,
environment placement raises preparation, and stop rule limits in-session
redesign cost.
\end{remark}
